Table~\ref{tab:vector:ge:activity} lists the commands that are available for activities for vectors of group elements.
\orbitertableofcommands{vector_ge_activity}


\bigskip

The following command computes and prints the conjugacy classes in $\PGGL(2,4)$:

\sourcecode{PGGL_2_4_classes_reps}

The class representatives are chosen by Magma.

\bigskip



\bigskip

The next command computes and prints 
the rational normal form for the representatives of the conjugacy classes in $\PGL(4,4)$

\sourcecode{PGL_4_4_classes_reps}

\bigskip




The next command computes and prints 
the conjugacy classes in $\PGGL(4,4)$:

\sourcecode{PGGL_4_4_classes_reps}

The class representatives are chosen by Magma.

\bigskip






The next command illustrates how to obtain the generators of a group as a vector of group elements object. 
In the example, we consider the group $\Sym(6).$ 

\sourcecode{Sym_6_generators}

\bigskip


The command produces a latex report of the generators, which are the Coxeter generators.

\orbiteroutput{GRAPHICS/LATEX/Sym_6_generators}


\bigskip


The next command illustrates how to obtain 
the elements of a group as a vector of group elements object. 
In the example, we consider the group 
$\Sym(6)$ with $720$ elements. 
The command produces a latex report of the group elements.

\sourcecode{Sym_6_elements}

\bigskip



The following command creates the conjugacy 
class of orthogonal reflection in $\PGO^+(6,2)$ and 
pushes them back into $\PGL(4,2)$ using the inverse of the exterior square action.
Because the group elements lie in $\PGO^+(6,2)$ but not 
in the subgroup isomorphic to $\PGL(4,2)$, 
a fixed polarity is divided off before the push back.
After that, a report of the 28 group elements is produced. 
Finally, the group elements are exported to GAP.
Using GAP, the order of the group generated 
by these elements turns out to be 20160.

\bigskip





The command 

\sourcecode{GL_2_3_RNF_elements}

computes the rational normal form of each element in $\GL(2,3)$.
A transformation matrix is computed as well. The output is:

\orbiteroutput{GRAPHICS/LATEX/GL_2_3_RNF_elements}




\bigskip




\sourcecode{PGL_4_2_orthogonal_reflections_report_vector_ge}

\bigskip


The command produces a latex report of the group elements, 
which are matrices of size 4. 
The associated permutation action 
on the 15 points of $\PG(3,2)$ is shown also.
For reasons of space, the output has been truncated.

\orbiteroutput{GRAPHICS/LATEX/PGL_4_2_orthogonal_reflections_cut}



\bigskip







